\section{Metodología Experimental}
Todos los entrenamientos fueron realizados en la misma computadora de forma secuencial utilizando Docker para simular un entorno multicomputadora, sumando un total de \num{10.5} horas de ejecución. Las especificaciones técnicas del entorno se detallan en la \Cref{tab:environment}, mientras que las variables lógicas y configuraciones del modelo O.N.O. se consolidan en la \Cref{tab:configs}.

\begin{table}[ht]
	\caption{Configuración de Ensayos}
	\label{tab:configs}
	\centering
	\footnotesize
	\begin{tabular}{llp{2.4cm}}
		\toprule
		\textbf{Categoría} & \textbf{Parámetro} & \textbf{Valores Evaluados} \\
		\midrule
		\multirow{3}{*}{Modelo y Datos} & Dataset & MNIST (Full)~\cite{lecun1998gradient} \\
		 & Arquitectura de Red & LeNet-5 \\
		 & Función de Pérdida & Entropía Cruzada \\
		\midrule
		\multirow{3}{*}{Optimización} & Optimizador & Adam \newline $\alpha = 0.01 \newline \beta_1 = 0.9 \newline \beta_2 = 0.999 \newline \epsilon = 10^{-8}$ \\
		 & Tamaño de Mini-Batch & 64 \\
		 & Épocas Máximas & 48 \\
		\midrule
		Infraestructura & Serialización & Sparse ($r = 0.5$) \\
		\midrule
		\multirow{3}{*}{Variables Libres} & Algoritmo Distribuido & All-Reduce, Parameter Server (Sincrónico/Bloqueante) \\
		 & Nodos & 2, 4, 6, 8, 10 \\
		 & Épocas Offline & 0, 1, 2, 3, 4 \\
		\bottomrule
	\end{tabular}
	\begin{tablenotes}
		\item[*] Para el algoritmo de \emph{Parameter Server}, el total de nodos se compone de forma que la mitad de ellos son \emph{workers} y la otra \emph{servers}.
	\end{tablenotes}
\end{table}

\subsection{Compresión de Red: Serialización Sparse}
La configuración de la sección de infraestructura responde al método de serialización de gradientes del sistema, el cual soporta una variante \emph{sparse} (o dispersa) diseñada para optimizar la comunicación distribuida limitando el intercambio de sectores poco impactantes de los gradientes~\cite{aji-heafield-2017-sparse}. El parámetro $r$ sirve para acotar el sampleo del gradiente que se quiere transmitir: a mayor $r$, mayor será la compresión teórica. 

El algoritmo calcula la cardinalidad efectiva del mensaje mediante:

\begin{equation}
k = \text{round}\left(|g| \cdot (1.0 - r)\right)
\end{equation}

El valor resultante $k$ se utiliza para escoger los $k$ elementos de mayor magnitud absoluta dentro del tensor de gradientes local, fijando un umbral mínimo. Todo valor que alcance o supere dicho umbral se incorporará al mensaje por enviar. De esta forma, un valor de $r = 0.0$ samplea la totalidad del gradiente, mientras que un $r = 0.5$ (como el configurado para estos experimentos) filtra el \SI{50}{\percent} de los componentes para disminuir el transito por red. Los elementos descartados se acumulan en un gradiente residual para ser considerados en el próximo mensaje.

\begin{table}[ht]
	\caption{Entorno de Ejecución}
	\label{tab:environment}
	\centering
	\footnotesize
	\begin{tabular}{llp{2.4cm}}
		\toprule
		\textbf{Categoría} & \textbf{Componente} & \textbf{Detalle} \\
		\midrule
		\multirow{5}{*}{Procesador} & Modelo & Intel Core i5-8350U \\
			& Arquitectura & x86\_64 \\
			& Núcleos Físicos & 4 \\
			& Hilos Lógicos & 8 \\
			& Frecuencia de Reloj & Base: \SI{1.7}{\giga\hertz} \newline Turbo: \SI{3.6}{\giga\hertz} \\
		\midrule
		\multirow{2}{*}{Memoria} & Capacidad Total RAM & 8 GB \\
			& Tipo & DDR4 \\
			& Velocidad & \SI{2400}{\mega\hertz} \\
		\midrule
		\multirow{2}{*}{Entorno de Software} & Sistema Operativo & Ubuntu 24.04 LTS \\
			& Entorno de Ejecución & Docker v29.6.1 \newline Rust v1.96.1 \newline CPython v3.14.0 \\
		\bottomrule
	\end{tabular}
\end{table}

\FloatBarrier

